\section{Evaluation}
\label{sec:evaluation}

Evaluation should match the component or system claim being made. A gain in final-answer accuracy may result from better retrieval, a stronger reader, additional context, memorized knowledge, or a different stopping policy. Multi-hop evaluation therefore requires both end-to-end outcomes and process-level measurements that expose where information was found, lost, ignored, or reused. The four functional categories in this survey provide a claim-aligned evaluation template: initial retrieval concerns the entry pool, follow-up retrieval concerns trajectories, evidence integration concerns preservation and use, and feedback learning concerns persistent improvement across tasks.

\subsection{Retrieval Evaluation}
\label{subsec:retrieval_evaluation}

Standard retrieval metrics include Recall@$K$, precision, mean reciprocal rank, and nDCG. For multi-hop tasks, these metrics should be computed over different evidence roles. \emph{Bridge recall} measures whether evidence needed to formulate a later query is available. \emph{Answer-evidence recall} measures whether answer-bearing sources are available. \emph{Complete-chain recall} measures whether all required units are present simultaneously. A system can have high document recall while failing complete-chain recall if it repeatedly retrieves the same hop.

Initial and follow-up retrieval should also be evaluated over budget curves. Reporting a single $K$ hides whether a method improves retrieval or simply requires a larger candidate pool. Useful curves vary candidate depth, number of retrieval rounds, total retrieved tokens, and unique sources. For graph systems, analogous quantities include subgraph size, path coverage, and expansion depth. When the dataset supplies supporting facts, as in HotpotQA and MuSiQue~\cite{yang2018hotpotqa,trivedi2022musique}, retrieval metrics should be computed before reader filtering so that retrieval and evidence use remain separable.

Candidate-pool control is essential for reranking and selection claims. Comparisons should fix the first-stage pool, context budget, and reader. Oracle analyses can estimate the headroom available if the correct evidence is selected, but oracle evidence should not be mixed with open-domain results. Datasets may also contain disconnected reasoning shortcuts; answer accuracy should therefore be interpreted together with evidence coverage and shortcut-resistant subsets~\cite{trivedi2020nohotpot}.

\subsection{Search Trajectory Evaluation}
\label{subsec:trajectory_evaluation}

Follow-up retrieval produces a sequence rather than a single ranking. Trajectory evaluation should measure whether each action advances an unresolved information need. For decomposition methods, relevant metrics include subquestion validity, dependency preservation, answerability, and coverage of the original task. For iterative retrieval, useful metrics include hop success, path recall, recovery after an incorrect hop, branching factor, and the fraction of trajectories that terminate with a complete support set.

Adaptive and agentic systems require action-level diagnostics. These include retrieval-trigger precision and recall, invalid-action rate, query duplication, source revisit rate, premature stopping, unnecessary retrieval, and the number of tool or browser actions. Cost should be reported in at least two units: environment interactions, such as searches and pages opened, and model consumption, such as tokens or inference calls. Wall-clock time and monetary API cost are useful but hardware- and provider-dependent.

Trajectory quality cannot be reduced to length. A shorter trajectory may be efficient or may stop before the task is solved; a longer one may be wasteful or may correctly verify conflicting evidence. Cost-adjusted metrics should therefore condition on task success, for example successful tasks per search call or answer quality under a fixed action budget. Releasing trajectories, query logs, and stopping decisions enables later analysis and should be considered part of the evaluation artifact when privacy and licensing permit.

\subsection{Evidence Integration Evaluation}
\label{subsec:integration_evaluation}

Evidence-integration evaluation asks whether acquired information survives selection and compression, is presented in a usable form, and causally supports the output. Membership and ordering must be separated. To test a selector, fix the candidate pool and presentation policy. To test ordering, fix the evidence set. To test compression, compare the compressed representation with the same uncompressed sources under matched reader conditions.

Selection metrics include complete-support preservation, redundancy, diversity, contradiction rate, and reader utility. Compression metrics include compression ratio, token and latency reduction, support retention, relation retention, and provenance retention. Context-construction studies should sweep context length and evidence position because long-context models can degrade as additional retrieved passages are added~\cite{liu2024lostmiddle,jin2025longcontextrag}.

Grounding requires claim-level checks. Citation correctness asks whether a cited source supports the associated claim; citation completeness asks whether claims that require evidence are cited. ALCE formalizes several aspects of citation-grounded generation~\cite{gao2023alce}. Supporting-fact or rationale overlap is useful when annotations exist, but generated rationales may be post hoc. Stronger tests intervene on evidence: remove a required source, replace a relation, permute dependent evidence, or introduce a contradictory source. If the output does not respond appropriately, the system may not be using the claimed evidence.

Long-form synthesis adds coverage, coherence, temporal consistency, and source-quality dimensions. Overall report scores should be accompanied by claim--source mappings and per-claim judgments. Human evaluators should have access to the cited sources; otherwise, fluency and topical plausibility can dominate the score.

\subsection{Feedback Learning Evaluation}
\label{subsec:feedback_evaluation}

Feedback-learning claims require a sequence of tasks or repeated rounds with persistent state. A successful retry on the same task is insufficient evidence of cross-task learning unless the retained update is isolated. Evaluation should distinguish \emph{within-task recovery}, \emph{same-distribution transfer}, and \emph{out-of-distribution transfer}.

For experience memories, report memory write rate, retrieval hit rate, retrieved-memory relevance, context overhead, and performance with memory retrieval disabled. For skill libraries, report skill selection accuracy, execution success, reuse frequency, library growth, routing cost, and composition failures. For learned retrievers or policies, report learning curves, held-out-task performance, reward definitions, and ablations that freeze each jointly evolving component.

Negative transfer and forgetting are first-class outcomes. A growing memory can retrieve stale or conflicting lessons; a revised skill can improve one task family while damaging another; and a retriever trained on agent-generated queries can overfit to the current policy. Evaluation should include chronological tests, memory-size sweeps, deletion or replacement interventions, and performance on tasks not represented in the retained experience. The strongest evidence links a specific retained object to a changed action and a downstream improvement on a later task.

\subsection{Benchmarks}
\label{subsec:benchmarks}

Benchmarks differ in which parts of the process they make observable. HotpotQA supplies supporting facts and diverse bridge and comparison questions~\cite{yang2018hotpotqa}. MuSiQue constructs questions from composed single-hop problems and reduces several disconnected shortcuts~\cite{trivedi2022musique}. HybridQA and MultiModalQA expand the evidence substrate to tables, text, and images~\cite{chen2020hybridqa,talmor2021multimodalqa}. These datasets are useful for controlled component evaluation because evidence sets or intermediate structure are partly annotated.

Open-web and deep-search evaluations trade annotation density for ecological validity. Browser-assisted systems such as WebGPT operate over changing web content and action interfaces~\cite{nakano2021webgpt}. Such benchmarks should snapshot sources or log timestamps, search-provider settings, and page content so that results can be reproduced. Long-form and citation-oriented benchmarks should separate retrieval, source use, and writing quality rather than collapsing them into one judge score.

No benchmark covers all four categories. A task with gold supporting facts may reveal retrieval and integration errors but provide no repeated-task signal for feedback learning. An agent benchmark may expose trajectories but lack evidence annotations. A continual-learning benchmark may measure transfer without source grounding. Table~\ref{tab:evaluation_alignment} therefore presents a minimum alignment between claims, metrics, and controls.

\begin{table*}[tbp]
\centering
\caption{Claim-aligned evaluation for multi-hop retrieval and reasoning.}
\label{tab:evaluation_alignment}
\small
\begin{tabularx}{\textwidth}{@{}p{2.8cm}p{4.1cm}Y Y@{}}
\toprule
Claim & Minimum direct measurements & Required controls & Common confound \\
\midrule
Better initial retrieval & Bridge recall, answer-evidence recall, complete-chain Recall@$K$ & Fixed corpus, candidate depth, and query & Stronger reader or larger candidate pool \\
Better follow-up retrieval & Hop/path success, query/action trace, stopping, total search cost & Fixed entry pool and action budget & More retrieval calls or hidden tool access \\
Better evidence integration & Support preservation, ordering/compression test, grounding or citation metrics & Fixed evidence membership, reader, and context budget & Selection and reader changes occur together \\
Better long-form synthesis & Claim coverage, citation correctness/completeness, contradiction handling & Source snapshot and claim-level evaluation & Fluency dominates judge score \\
Better feedback learning & Learning curve, memory/skill use, held-out transfer, negative transfer & Persistent-state ablation and matched context cost & Repeated sampling or within-task retry \\
Better end-to-end system & Answer/report quality plus retrieval, trajectory, grounding, and cost & Same data access, tools, model budget, and stopping limits & Unreported differences in resources or model scale \\
\bottomrule
\end{tabularx}
\end{table*}

\paragraph{Diagnostic audit.}
The companion audit uses a stratified \AuditSize-paper selection across the four functional categories. Each record is coded only after a full-text source passage, table, figure, or intervention is located. The codebook distinguishes \emph{reported}, \emph{not reported}, \emph{not applicable}, and \emph{unclear}; absence is never inferred from a search snippet or abstract. This design allows the survey to summarize evaluation practice without treating the larger discovery corpus as fully audited evidence.

\paragraph{Summary.}
Multi-hop evaluation is credible when it reveals both outcome and process. The required measurements depend on the claim, but the general rule is stable: fix neighboring components, report budgets and trajectories, and test whether the evidence or experience attributed to success is actually used.
